%! Mean, Variance, and Covariance — Unit 8, Lesson 8.4 — Student Packet Cover Sheet
\documentclass[10pt]{article}
\usepackage{linalg-article}
\usepackage{linalg-boxes}
\usepackage{ltablex}
\keepXColumns

\begin{document}

% Full-bleed burgundy banner
\begin{tikzpicture}[remember picture, overlay]
  \fill[burgundy] (current page.north west)
    rectangle ([yshift=-0.9in]current page.north east);
\end{tikzpicture}
\vspace{-0.8in}

\begin{center}
  {\color{white}\textbf{\LARGE Linear Algebra}}\\[4pt]
  {\color{blushmid}\large\textbf{Unit 8 \quad Learning from Data}}\\[6pt]
  {\color{white}\textbf{\large Lesson 8.4 \quad Mean, Variance, and Covariance}}
\end{center}

\vspace{0.22in}
\namedateperiod
\vspace{0.18in}

\begin{learningtargetbox}
\small
\begin{enumerate}[leftmargin=1.5em, itemsep=5pt]
  \item \ldots{} find the \textbf{mean} (center) and \textbf{variance} (spread, the average squared distance from the mean) of a data set.
  \item \ldots{} compute the \textbf{covariance} of two features and read its \textbf{sign} --- do they move together, oppositely, or not at all?
  \item \ldots{} build the symmetric \textbf{covariance matrix} $V$ and explain why its eigenvectors are the \textbf{principal components} (Unit~7).
\end{enumerate}
\end{learningtargetbox}

\vspace{0.14in}

\begin{tocbox}
\small
\renewcommand{\arraystretch}{1.5}
\begin{tabularx}{\linewidth}{@{} c l X r @{}}
\rowcolor{blush}
\textbf{\#} & \textbf{Component} & \textbf{Description} & \textbf{Score} \\
\midrule
1 & Warm-Up        & Spiral review: an average, a sum of squared distances, and eigenvalues of a symmetric matrix & \blank{1.2cm} \\
2 & Guided Notes   & Mean $=$ center; variance $=$ spread; covariance $=$ moving together; the covariance matrix $V$ & \blank{1.2cm} \\
3 & Group Activity & Spread of one feature, the covariance matrix of two features, and its principal components & \blank{1.2cm} \\
4 & Exit Ticket    & Quick check: a mean \& variance, reading a covariance matrix, and why $V$ is symmetric & \blank{1.2cm} \\
5 & Homework       & Variance \& spread, covariance \& its sign, building $V$, the PCA link, and justifications & \blank{1.2cm} \\
\midrule
  & \multicolumn{2}{r}{\textbf{Total}} & \blank{1.2cm} \\
\end{tabularx}
\end{tocbox}

\vspace{0.10in}

\begin{remindbox}
\small A data set is a \textbf{cloud of points}. Three numbers describe its shape: the \textbf{mean} $\mu$
(where it centers), the \textbf{variance} $\sigma^2=\frac1N\sum(x_i-\mu)^2$ (how spread out it is), and the
\textbf{covariance} $\sigma_{xy}=\frac1N\sum(x_i-\mu_x)(y_i-\mu_y)$ (how two features move together --- $+$
together, $-$ oppositely, $0$ unrelated). Packed into the symmetric \textbf{covariance matrix}
$V=\left[\begin{smallmatrix}\sigma_x^2 & \sigma_{xy}\\ \sigma_{xy} & \sigma_y^2\end{smallmatrix}\right]$, these
are exactly what PCA (7.3) works with: $V$'s eigenvectors are the \textbf{principal components}, the directions
of greatest spread. Understanding data is understanding a symmetric matrix --- the last idea of the course.
\end{remindbox}

\end{document}
